\documentclass{article}
\usepackage{amsmath}
\begin{document}

\section{Hashmark LaTeX Preview}

This is a quick test of the new \textbf{side-by-side LaTeX preview} in Hashmark.
You can write LaTeX on the left and watch it render on the right with a 300ms debounce.

\subsection{Math support via KaTeX}

Inline math: $E = mc^2$ and $\sum_{i=1}^{n} i = \frac{n(n+1)}{2}$.

Display math:
\[
\int_0^\infty e^{-x^2}\, dx = \frac{\sqrt{\pi}}{2}
\]

Aligned equations:
\begin{align*}
(a+b)^2 &= a^2 + 2ab + b^2 \\
(a-b)^2 &= a^2 - 2ab + b^2
\end{align*}

\subsection{Lists}

Items I like:
\begin{itemize}
  \item \emph{Markdown} for writing
  \item \textbf{LaTeX} for math-heavy notes
  \item Nested:
  \begin{itemize}
    \item One
    \item Two
  \end{itemize}
\end{itemize}

Numbered:
\begin{enumerate}
  \item First
  \item Second with code \texttt{like this}
  \item Third
\end{enumerate}

\subsection{Links and quotes}

Visit \href{https://example.com}{this link} or \url{https://hashmark.app}.

\paragraph{A paragraph head} This is a paragraph-titled section with inline math $a^2 + b^2 = c^2$.

\subsection{Edge cases}

Special chars: 5 < 7, true \& false, 50\% off, \$10, snake\_case, \#hashtag.

Unknown macro: \customcommand{this content still shows}.

\cite{einstein1905} A reference placeholder.

\end{document}
